% Part: lambda-calculus
% Chapter: syntax
% Section: terms

\documentclass[../../../include/open-logic-section]{subfiles}

\begin{document}

\olfileid{lam}{syn}{trm}
\olsection{Suku}

Suku-suku kalkulus lambda dibangun kanthi
induktif saka variabel tanpa wates cacahé
$\Obj{v_0}$, $\Obj{v_1}$, \dots, lambang
``$\lambd$'', lan tandha kurung. Kita bakal
nggunakake $x$, $y$, $z$, \dots{} kanggo
nyebut variabel, lan $M$, $N$, $P$, \dots{}
kanggo nyebut suku.

\begin{defn}[Suku] \ollabel{defn:term}
Himpunan \emph{suku} kalkulus lambda
ditetepake kanthi induktif kaya mangkene:
\begin{enumerate}
  \item \ollabel{defn:term-var} Yen $x$ iku
    variabel, mula $x$ iku suku.
  \item \ollabel{defn:term-abs} Yen $x$ iku
    variabel lan $M$ iku suku, mula
    $(\lambd[x][M])$ iku suku.
  \item \ollabel{defn:term-app} Yen $M$
    lan $N$ iku suku, mula $(MN)$ iku suku.
\end{enumerate}
\end{defn}

Yen suku $(\lambd[x][M])$ dibentuk miturut
\olref{defn:term-abs}, suku iku diarani
asil \emph{abstraksi}, lan $x$ ing
$\lambd[x]$ diarani \emph{!!{parameter}}.
Suku $(MN)$ kang dibentuk miturut
\olref{defn:term-app} iku asil
\emph{aplikasi}.

Suku-suku kang nembe ditetepake nganggo
tandha kurung kanthi lengkap. Iki bisa
dadi rada ruwet, kaya kang katon ing
suku $(\lambd[x][((\lambd[x][x])(\lambd[x][(xx)]))])$.
Kita bakal ngenalake konvensi supaya
sawetara tandha kurung bisa disingkirake.
Nanging definisi resmi nggampangake kita
nemtokake carane suku dibentuk miturut
\olref{defn:term}. Umpamane, langkah
pungkasan nalika mbentuk suku
$(\lambd[x][((\lambd[x][x])(\lambd[x][(xx)]))])$
mesthi abstraksi kanthi !!{parameter}~$x$.
Suku iki diasilake kanthi abstraksi saka
suku $((\lambd[x][x])(\lambd[x][(xx)]))$,
yaiku aplikasi saka rong suku. Saben
suku ing pasangan iku asil abstraksi,
lan sateruse.

\begin{prob}
% OLPL-241: Kurung aplikasi njaba diwenehake supaya manut definisi resmi.
Andharna carane
$(\lambd[g][((\lambd[x][(g (x x))])
  (\lambd[x][(g (x x))]))])$ dibentuk.
\end{prob}

\end{document}
